% Unnumbered chapter (no large chapter number in the heading), in TOC. Do not
% \refstepcounter{chapter} so the following chapter keeps the next number (no
% gap in the sequence). book's \chapter starts with \cleardoublepage/\clearpage;
% relax both so the heading can follow the previous text on the same page.
\begingroup
\let\cleardoublepage\relax
\let\clearpage\relax
\vspace{25pt}
\chapter*{Barcelona as a case study}
\phantomsection
\label{ch:use_case_barcelona}
\addcontentsline{toc}{chapter}{\protect\numberline{}Barcelona as a case study}
\endgroup
\vspace{-45pt} % remove the extra space added by the chapter title

% \chapter* does not advance the chapter counter, so
% floats would otherwise inherit the previous chapter's number. Use Part III prefix.
\begingroup
\setcounter{figure}{0}
\setcounter{table}{0}
\renewcommand{\thefigure}{III.\arabic{figure}}
\renewcommand{\thetable}{III.\arabic{table}}

Part~\ref{pt:data-driven-station-planning} develops a flexible data-driven framework for station planning in docked \acp{bss}. Before presenting the methodological details in Chapter~\ref{ch:station_planning_in_barcelona}, this Section introduces Barcelona as it is the empirical setting in which the framework is tested. This Section is adapted from Grau-Escolano et al.~\cite{GrauEscolano2025}, published in \textit{Journal of Urban Mobility}.

Barcelona, the capital of Catalonia in northeastern Spain, occupies approximately 101 km² and was home to roughly 1.7 million residents in 2024, making it one of Europe’s most densely populated cities. Its municipal boundary is delineated by two important geographic features: the Mediterranean Sea to the southeast and the Collserola mountain range to the northwest. These natural limits not only constrain urban expansion but also shape travel patterns, as coastal flatlands yield to the steep slopes of Collserola rising to nearly 450 m (Figure~\ref{fig:Barcelona_use_case}C). Administratively, Barcelona is divided into 10 districts and 73 neighborhoods (Figure~\ref{fig:Barcelona_use_case}A), which display marked socioeconomic contrasts (Figure~\ref{fig:Barcelona_use_case}B and D). Districts such as Sarrià–Sant Gervasi and Les Corts stand out with the highest average incomes, while Nou Barris and parts of Ciutat Vella record much lower income levels. The l'Eixample district, although more socioeconomically mixed, concentrates the largest number of amenities and \acp{poi}, making it a major hub of urban activity. 


\begin{figure}[!htbp]
    \centering
    \includegraphics[width=0.85\textwidth]{0_plots/paper1/main/barcelona_six_panel_comparison.jpg}
    \caption[Overview of Barcelona's urban context]{
        \textbf{Overview of Barcelona's urban context.} 
        \textbf{(A)} Administrative division of Barcelona: 10 districts (colored polygons) and 73 neighborhoods (gray boundaries). 
        \textbf{(B)} Average net income per capita by neighborhood in 2022. 
        \textbf{(C)} Topographic contours showing elevation variations across the city.
        \textbf{(D)} Number of \acp{poi} per neighborhood. Obtained from \ac{osm} in 2025.
        \textbf{(E)} Bike lanes (blue lines) in 2025 overlaid on bike lane density by neighborhood.
        \textbf{(F)} \textit{Bicing} stations (red points) in 2025 overlaid on station density by neighborhood.
        }
    \label{fig:Barcelona_use_case}
\end{figure}

Over the past two decades, Barcelona has actively pursued policies to promote sustainable and active transport. Initiatives such as the superblock program, low-emission zones, and the expansion of the cycling network reflect this agenda. Today, the city offers more than 260 km of dedicated bike lanes and nearly 2,000 km of bike-friendly streets (Figure~\ref{fig:Barcelona_use_case}E), complemented by 8 metro lines, 6 tram lines, and over 100 bus routes with more than 8,000 stops. Beyond the city limits, rail networks (\textit{FGC} and \textit{Rodalies}) and interurban buses extend connectivity to other municipalities.

How residents make use of this infrastructure is reflected in the 2024 survey \textit{Enquesta de Mobilitat en Dia Feiner} (EMEF). According to this survey~\cite{EMEF2024}, Barcelona residents made 5.75 million weekday trips in 2024, excluding intensive work-related travel, which equals an average of 3.9 trips per person per day. Overall, active mobility is the leading mode, followed by \ac{pt} and private vehicles (Table~\ref{tab:modal_split}). However, patterns differ depending on the type of trip. For internal trips within Barcelona, active mobility dominates, whereas for trips connecting Barcelona with surrounding municipalities, \ac{pt} is the main mode. Within active mobility, the vast majority of trips are made on foot, while cycling and personal mobility vehicles (e.g., e-scooters, segways, or similar devices) together account for less than 4\% of trips.

\begin{table}[!htbp]
    \centering
    \caption[Modal split of trips in Barcelona based on EMEF 2024 survey]{
        \textbf{Modal split of trips in Barcelona based on EMEF 2024 survey~\cite{EMEF2024}.}
    }
    \label{tab:modal_split}
    \begin{tabular}{lccc}
        \hline
        & \textbf{Total (\%)} & \textbf{Internal (\%)} & \textbf{Connections (\%)} \\
        \hline
        Active mobility & 54.5 & 60 & 5 \\
        \quad Walking & 51.3 & -- & -- \\
        \quad Bicycle & 2.4 & -- & -- \\
        \quad E-scooters, segways, etc. & 0.8 & -- & -- \\
        Public transport & 28.8 & 28 & 53 \\
        Private vehicles & 16.7 & 12 & 42 \\
        \hline
    \end{tabular}
\end{table}

Within this context, \textit{Bicing} has grown since its 2007 launch into a system of 557 stations (Figure~\ref{fig:Barcelona_use_case}F) and nearly 8,000 bicycles as of 2025, of which 5,000 are \acp{e-b}. This supply growth has been accompanied by rising demand: by the end of 2023, the system had approximately 147,000 subscribers and recorded 18 million trips, according to Barcelona's City Council~\cite{AjuntamentBarcelona2024_Bicing}. Empirical studies also reveal distinct usage patterns~\cite{Grau-Escolano2024}. \acp{e-b} are used more intensively in hilly districts near Collserola, whereas \acp{m-b} dominate in the city's flatter coastal areas. A more detailed description of the system, including its design and pricing, is given in \hyperref[ch:use_case_bicing]{\textit{Bicing as a case study}} at the beginning of Part~\ref{pt:predictive_operations}, and how citizens use it is analysed in Chapter~\ref{ch:bicing_mobility_patterns}.

These characteristics make Barcelona a compelling case study. The city combines compactness and high population density with marked topographic contrasts between flat coastal areas and hilly districts, as well as pronounced socioeconomic diversity. Although active mobility is widespread, cycling remains marginal at only 2.4\% of trips, highlighting clear potential for growth. This aligns with the city's Urban Mobility Plan 2025--2030, aims to increase the combined share of cycling and personal mobility vehicles by nearly 50\% by 2030~\cite{PMU2025}. Together with the availability of extensive open data, these conditions make Barcelona a suitable real-world setting for evaluating the proposed framework.


\endgroup